% GENERATED FILE. Edit canonical lessons, then run npm run generate:books.
% canonical-source-hash: fnv1a64:3c353bd152cfd8d3
% canonical-lessons: GE-C489-verwandt, GE-C489-verwandtschaft, GE-C489-neffe, GE-C489-nichte, GE-C489-schwager, GE-R489-review-482-489

\chapter{Neffen und Nichten}
\label{ch:ge-neffen-und-nichten}
\hlchaptermodality{489}

\begin{chapteropening}
\textbf{By the end of this chapter:} \emph{I can talk about my wider family.}

Complete the last lesson of chapter 489: I can talk about my wider family.
\end{chapteropening}

\section[verwandt]{verwandt — related (by family)}

\label{lesson:GE-C489-verwandt}



\begin{quote}
Before the new one: say the German for to get some rest, then the German for to recover.
\end{quote}

\subsection*{You'll want to know: verwandt}
\textbf{verwandt} — \textquotedblleft{}related (by family)\textquotedblright{}.

\begin{quote}
Wir haben den gleichen Nachnamen, aber wir sind nicht verwandt.
\end{quote}

\begin{quote}
\emph{We have the same surname, but we are not related.}
\end{quote}

\emph{Verwandt mit} + dative. \emph{Der Verwandte} or \emph{die Verwandte} — a relative.

\begin{cousinweb}
\emph{verwandt} is the old past participle of \emph{verwenden} \textquotedblleft{}to turn, to turn towards\textquotedblright{} (EtymDB), built on \emph{wenden} \textquotedblleft{}to turn\textquotedblright{}.
\end{cousinweb}

\begin{grammarlens}[title={What you've built}]
One more word for telling, explaining and describing at B1.
\end{grammarlens}

\subsection*{Your turn}
\begin{itemize}
\raggedright
  \item \emph{Say it:} \emph{verwandt}
  \item \emph{Say it:} \emph{verwandt}, once more
  \item \emph{Say it:} a sentence of your own with \emph{verwandt}
  \item \emph{Recall:} say the German for to get some rest, then the German for to recover, then say \emph{verwandt} again
\end{itemize}

\begin{tcolorbox}[breakable,colback=teal!4,colframe=teal!35!black,title={Before you move on}]
What does verwandt mean? (\textquotedblleft{}Related (by family)\textquotedblright{}.) Say it once more.
\end{tcolorbox}
\section[die Verwandtschaft]{die Verwandtschaft — relatives; kinship}

\label{lesson:GE-C489-verwandtschaft}



\begin{quote}
Before the new one: say the German for to recover, then the German for related.
\end{quote}

\subsection*{You'll want to know: die Verwandtschaft}
\textbf{die Verwandtschaft} — \textquotedblleft{}relatives; kinship\textquotedblright{}.

\begin{quote}
Zu Weihnachten kommt die ganze Verwandtschaft.
\end{quote}

\begin{quote}
\emph{All the relatives come for Christmas.}
\end{quote}

A collective noun: all your relatives together. Plural, rare: \emph{die Verwandtschaften}.

\begin{cousinweb}
\emph{Verwandtschaft} is built on \emph{verwandt} \textquotedblleft{}related\textquotedblright{} with \emph{-schaft}; \emph{verwandt} is the old participle of \emph{verwenden} \textquotedblleft{}to turn towards\textquotedblright{} (EtymDB).
\end{cousinweb}

\begin{grammarlens}[title={What you've built}]
One more word for telling, explaining and describing at B1.
\end{grammarlens}

\subsection*{Your turn}
\begin{itemize}
\raggedright
  \item \emph{Say it:} \emph{die Verwandtschaft}
  \item \emph{Say it:} \emph{die Verwandtschaft}, once more
  \item \emph{Say it:} a sentence of your own with \emph{die Verwandtschaft}
  \item \emph{Recall:} say the German for to recover, then the German for related, then say \emph{die Verwandtschaft} again
\end{itemize}

\begin{tcolorbox}[breakable,colback=teal!4,colframe=teal!35!black,title={Before you move on}]
What does die Verwandtschaft mean? (\textquotedblleft{}Relatives; kinship\textquotedblright{}.) Say it once more.
\end{tcolorbox}
\section[der Neffe]{der Neffe — nephew}

\label{lesson:GE-C489-neffe}



\begin{quote}
Before the new one: say the German for related, then the German for relatives.
\end{quote}

\subsection*{You'll want to know: der Neffe}
\textbf{der Neffe} — \textquotedblleft{}nephew\textquotedblright{}.

\begin{quote}
Mein Neffe ist zehn und spielt gern Fußball.
\end{quote}

\begin{quote}
\emph{My nephew is ten and likes playing football.}
\end{quote}

Weak noun: \emph{den Neffen}, \emph{dem Neffen}. Plural: \emph{die Neffen}.

\begin{cousinweb}
\emph{Neffe} is inherited from Old High German \emph{nefo} (EtymDB), an old Germanic family word.
\end{cousinweb}

\begin{grammarlens}[title={What you've built}]
One more word for telling, explaining and describing at B1.
\end{grammarlens}

\subsection*{Your turn}
\begin{itemize}
\raggedright
  \item \emph{Say it:} \emph{der Neffe}
  \item \emph{Say it:} \emph{der Neffe}, once more
  \item \emph{Say it:} a sentence of your own with \emph{der Neffe}
  \item \emph{Recall:} say the German for related, then the German for relatives, then say \emph{der Neffe} again
\end{itemize}

\begin{tcolorbox}[breakable,colback=teal!4,colframe=teal!35!black,title={Before you move on}]
What does der Neffe mean? (\textquotedblleft{}Nephew\textquotedblright{}.) Say it once more.
\end{tcolorbox}
\section[die Nichte]{die Nichte — niece}

\label{lesson:GE-C489-nichte}



\begin{quote}
Before the new one: say the German for relatives, then the German for nephew.
\end{quote}

\subsection*{You'll want to know: die Nichte}
\textbf{die Nichte} — \textquotedblleft{}niece\textquotedblright{}.

\begin{quote}
Meine Nichte hat mir ein Bild gemalt.
\end{quote}

\begin{quote}
\emph{My niece painted me a picture.}
\end{quote}

Plural: \emph{die Nichten}.

\begin{cousinweb}
\emph{Nichte} was borrowed from Middle Low German \emph{nichte} (EtymDB), the northern form of an old Germanic word for a niece.
\end{cousinweb}

\begin{grammarlens}[title={What you've built}]
One more word for telling, explaining and describing at B1.
\end{grammarlens}

\subsection*{Your turn}
\begin{itemize}
\raggedright
  \item \emph{Say it:} \emph{die Nichte}
  \item \emph{Say it:} \emph{die Nichte}, once more
  \item \emph{Say it:} a sentence of your own with \emph{die Nichte}
  \item \emph{Recall:} say the German for relatives, then the German for nephew, then say \emph{die Nichte} again
\end{itemize}

\begin{tcolorbox}[breakable,colback=teal!4,colframe=teal!35!black,title={Before you move on}]
What does die Nichte mean? (\textquotedblleft{}Niece\textquotedblright{}.) Say it once more.
\end{tcolorbox}
\section[der Schwager]{der Schwager — brother-in-law}

\label{lesson:GE-C489-schwager}



\begin{quote}
Before the new one: say the German for nephew, then the German for niece.
\end{quote}

\subsection*{You'll want to know: der Schwager}
\textbf{der Schwager} — \textquotedblleft{}brother-in-law\textquotedblright{}.

\begin{quote}
Mein Schwager hilft uns beim Umzug.
\end{quote}

\begin{quote}
\emph{My brother-in-law is helping us move.}
\end{quote}

Plural: \emph{die Schwager}. A sister-in-law is \emph{die Schwägerin}.

\begin{cousinweb}
\emph{Schwager} is inherited from Middle High German \emph{swāger} and Old High German \emph{suāgur}, from Proto-Germanic \emph{swēgraz} \textquotedblleft{}husband's brother\textquotedblright{} (EtymDB).
\end{cousinweb}

\begin{grammarlens}[title={What you've built}]
One more word for telling, explaining and describing at B1.
\end{grammarlens}

\subsection*{Your turn}
\begin{itemize}
\raggedright
  \item \emph{Say it:} \emph{der Schwager}
  \item \emph{Say it:} \emph{der Schwager}, once more
  \item \emph{Say it:} a sentence of your own with \emph{der Schwager}
  \item \emph{Recall:} say the German for nephew, then the German for niece, then say \emph{der Schwager} again
\end{itemize}

\begin{tcolorbox}[breakable,colback=teal!4,colframe=teal!35!black,title={Before you move on}]
What does der Schwager mean? (\textquotedblleft{}Brother-in-law\textquotedblright{}.) Say it once more.
\end{tcolorbox}
\section[(dialogue)]{Words from chapters 482–489 — a review}

\label{lesson:GE-R489-review-482-489}



\begin{quote}
Say the German for to sprain and bruise.
\end{quote}

\subsection*{Your turn}
Say these aloud:

\begin{itemize}
\raggedright
  \item to sprain, bruise, bandage, to disinfect, scar
  \item plaster cast, crutch, to limp, wheelchair, to heal
  \item forehead, cheek, chin, eyebrow, eyelash
  \item back of the neck, elbow, wrist, ankle, hip
  \item tiredness, exhausted, to yawn, to snore, nightmare
  \item family doctor, specialist, to examine, to treat, to prescribe
  \item to infect, infection, to vaccinate, to get some rest, to recover
  \item related, relatives, nephew, niece, brother-in-law
\end{itemize}

\begin{tcolorbox}[breakable,colback=teal!4,colframe=teal!35!black,title={Before you move on}]
To sprain? (\textbf{verstauchen}.) Bruise? (\textbf{die Prellung}.) Bandage? (\textbf{der Verband}.) To disinfect? (\textbf{desinfizieren}.) Scar? (\textbf{die Narbe}.) Plaster cast? (\textbf{der Gips}.) Crutch? (\textbf{die Krücke}.) To limp? (\textbf{humpeln}.) Wheelchair? (\textbf{der Rollstuhl}.) To heal? (\textbf{heilen}.) Forehead? (\textbf{die Stirn}.) Cheek? (\textbf{die Wange}.) Chin? (\textbf{das Kinn}.) Eyebrow? (\textbf{die Augenbraue}.) Eyelash? (\textbf{die Wimper}.) Back of the neck? (\textbf{der Nacken}.) Elbow? (\textbf{der Ellbogen}.) Wrist? (\textbf{das Handgelenk}.) Ankle? (\textbf{der Knöchel}.) Hip? (\textbf{die Hüfte}.) Tiredness? (\textbf{die Müdigkeit}.) Exhausted? (\textbf{erschöpft}.) To yawn? (\textbf{gähnen}.) To snore? (\textbf{schnarchen}.) Nightmare? (\textbf{der Albtraum}.) Family doctor? (\textbf{der Hausarzt}.) Specialist? (\textbf{der Facharzt}.) To examine? (\textbf{untersuchen}.) To treat? (\textbf{behandeln}.) To prescribe? (\textbf{verschreiben}.) To infect? (\textbf{anstecken}.) Infection? (\textbf{die Infektion}.) To vaccinate? (\textbf{impfen}.) To get some rest? (\textbf{erholen}.) To recover? (\textbf{genesen}.) Related? (\textbf{verwandt}.) Relatives? (\textbf{die Verwandtschaft}.) Nephew? (\textbf{der Neffe}.) Niece? (\textbf{die Nichte}.) Brother-in-law? (\textbf{der Schwager}.)
\end{tcolorbox}
